\documentclass{article}
\usepackage[UTF8]{ctex}
\usepackage{amsmath,amssymb,amsfonts}
\usepackage{enumitem}
\usepackage{geometry}
\geometry{a4paper, margin=1in}
% === 错题本排版优化（不参与题目标识） ===
\setlength{\parindent}{0pt}
\setlength{\parskip}{0.35em plus 0.10em minus 0.08em}
\setlength{\emergencystretch}{3em}
\setlist{itemsep=0.25em, topsep=0.40em, parsep=0pt}
\allowdisplaybreaks[2]
% === 错题本排版优化结束 ===
\begin{document}
\section*{东北大学2026年数学分析真题}
\section*{74．东北大学2026年数学分析试卷}

\subsection{题目1}
1．（15 分）求极限 \(\lim _{n \rightarrow \infty}\left[\ln \left(1+\frac{1}{n}\right)^{n}\right]^{n}\) ．

\subsection{题目2}
2．（15 分）将 \(f(x)=\ln x\) 按 \(\frac{x-1}{x+1}\) 进行幂级数展开．

\subsection{题目3}
3．（15 分）设 \(m, n \in \mathbb{N}^{+}\)，求 \(I(m, n)=\int_{0}^{1}(\ln x)^{m} x^{n} \mathrm{~d} x\) ．

\subsection{题目4}
4．（15 分）计算三重积分 \(\iiint_{\Omega} e^{|z|} \mathrm{d} x \mathrm{~d} y \mathrm{~d} z\) ，其中 \(\Omega\) 是单位球。

\subsection{题目5}
5．（15 分）叙述定义：\(x \rightarrow+\infty\) 时 \(f(x) \rightarrow-\infty ; x \rightarrow 0^{+}\)时，\(f(x)\) 不以 \(A\) 为极限．

\subsection{题目6}
6．（15 分）证明级数 \(\sum_{n=1}^{\infty}\left[n x e^{-n x}-(n-1) x e^{-(n-1) x}\right]\) 在 \([0,1]\) 上收敛但不一致收敛，而和函数连续．

\subsection{题目7}
7．（15 分）设 \(f(x)\) 在 \(\left[x_{1}, x_{2}\right]\) 上连续 \(\left(x_{1} x_{2}>0\right)\) ，在 \(\left(x_{1}, x_{2}\right)\) 上可导，证明：存在 \(\xi \in\left(x_{1}, x_{2}\right)\) ，使得

\[
\frac{1}{x_{2}-x_{1}}\left|\begin{array}{ll}
x_{1} & f\left(x_{1}\right) \\
x_{2} & f\left(x_{2}\right)
\end{array}\right|=\xi f^{\prime}(\xi)-f(\xi) .
\]

\subsection{题目8}
8．（15 分）设 \(f(x)\) 在 \((-\infty,+\infty)\) 上二阶连续可微，且 \(f(x)+f^{\prime \prime}(x)=-x g(x) f^{\prime}(x)\) ，其中 \(g(x)\) 为非负函数．证明：\(f(x)\) 有界．

\subsection{题目9}
9．（15 分）设 \(f(x)\) 为 \([a, b]\) 上的二阶可导函数，且 \(f^{\prime \prime}(x) \leq 0, f(x) \geq 0\) ，证明：

\[
\max _{x \in[a, b]} f(x) \leq \frac{2}{b-a} \int_{a}^{b} f(x) \mathrm{d} x
\]

\subsection{题目10}
10．（15 分）证明：\(f(x)=\left\{\begin{array}{l}\frac{1}{n}, \quad x=\frac{m}{n}(m, n \text { 是互质正整数 }) \\ 0, \quad x \text { 为无理数．}\end{array}\right.\) 在无理点连续，在有理点不连续．
\end{document}
